\section{Outro}\label{sec:outro}

What's next for dynamic stabilizer codes and Pauli webs?
Let's start with the former.
As touched on already in \Cref{ch:intro}, one can certainly argue that dynamic stabilizer codes have so far failed to live up to the hype that accompanied their arrival in \Ccite{Hastings_2021}.
They haven't got around no-go theorems that affect stabilizer codes~\cite{Fu_2025}, and in simulations they've only managed to outperform the surface code in architectures that allow for direct weight-two Pauli measurements (no auxiliary qubits needed)~\cite{dessertaine2025enhancedfaulttolerancephotonicquantum,Higgott_2024,Gidney_2021,Gidney_2022}.

On the other hand, it's early days for these codes, and we're still exploring some of the perks they offer.
For example, at any given time in a dynamic stabilizer code we have the freedom to choose any Pauli operator whatsoever to measure at the next timestep.
So for instance, using only a limited amount of this freedom, we can choose to repeat rounds of measurements multiple times in order to gain more information about the types of errors these measurements protect against.
In work inspired by \Cref{Higgott_2021} that hasn't made it into this thesis~\cite{derks2025dynamicalcodeshardwarenoisy}, we looked into whether exactly this idea could be used to tailor a certain class of dynamic stabilizer codes to biased noise models (where, say, $Z$ errors are more likely than other Pauli errors, or measurement errors are more likely than all Pauli errors).
Once again, the results were mostly underwhelming -- only in very limited regimes was this effective.
But this freedom could be pushed to its extreme -- rather than start with an operation schedule and repeat certain measurements, we could start with an effectively blank canvas and at each timestep adaptively choose what to measure, based on some function that balances gaining information about specific errors we may have accrued and maintaining the distance of the code.

All that said, I personally am pessimistic about the chances of these ideas bearing fruit.
I suspect instead that the lasting legacy of dynamic stabilizer codes will more high-level; they expanded our definition of what a stabilizer code can be, and led us to the realisation that in fact any Clifford circuit with Pauli measurements can already be viewed as a QEC code.
And this mindset shift has in turn had practical benefits, which we discuss more in the outro of the next chapter.

Now for Pauli webs.
So far, these have shown themselves to be very useful for reasoning about qubit Clifford ZX-diagrams.
A major open question for me is whether they can help us say anything about non-Clifford ZX-diagrams, and therefore about non-Clifford logical computation in QEC codes, which is necessary for universal quantum computing.
In general, going to the non-Clifford fragment of the ZX-calculus means a lot of tools break down, but there's been some nice work showing that you can still make a lot of progress if you restrict yourself to diagrams with \textit{dyadic} angles (those of the form $\frac{\pi}{2^k}$ for $k \in \mathbb{Z}$)~\cite{Kissinger_2024,KissingerWetering2024Book}.

To this end, something I would love to see explored here is the connection between Pauli webs, the ZH-calculus~\cite{Backens_2019} and the XP formalism~\cite{Webster_2022}.
The XP formalism is just the stabilizer formalism but for stabilizer groups that aren't necessarily subgroups of the Pauli group, but rather are groups generated by $X$ and $Z(\frac{\pi}{2^k})$ gates, for some fixed integer $k$ (hence the connection to ZX-diagrams with dyadic angles!).
Here we lose certain things like the efficient classical simulability of the Pauli stabilizer formalism, but this should be seen as a feature, not a bug, because anything classically simulable almost certainly can't reason about universal quantum computation.
In particular, the XP formalism can be used to reason about non-Clifford gates, e.g.\ as in \Ccite{davydova2025universalfaulttolerantquantum}.
Since Pauli webs can just be viewed as graphical Pauli stabilizer formalism, are there `XP webs' can be viewed as graphical XP stabilizer formalism?

The connection the ZH-calculus comes from the fact that XP stabilizer states are equivalent to \textit{hypergraph states}, which can be very naturally represented in the ZH-calculus.
This calculus can be viewed as an extension of the ZX-calculus, in which the two-legged Hadamard box is upgraded to a spider-like object with arbitrary many legs and a label $a \in \mathbb{C}$.
Of course, since the ZX-calculus is already complete, the ZH-calculus isn't actually any more expressive than the ZX-calculus, but certain diagrams (hypergraph states, for example!) are more natural in ZH than in ZX.

Another direction in which Pauli webs can be extended is to higher dimensional qudits and qumodes -- the latter being the subject of the second half of this thesis.
